\section{Zusammenfaßung}
Durch dieses Projekt haben wir die Möglichkeit gehabt, unser Wissen und Umgang mit Software Engineering in die Praxis einzusetzen. Während der Vorlesungszeiten haben wir uns mit wesentliche Themen beschäftigt, die uns geholfen hatten unsere Webanwendung zu programmieren sowie strukturiert die Anforderungen der Webanwendung zu bestimmen. Durch Persona, Anforderungsanalyse, use-case sowie Aktivitätsdiagramme haben wir die wichtgere Anforderung der Anwendung dargestellt. Durch Kombination von Servlet, MariaDB und Redis könnten wir SQL-Abfrage einfach gestalten sowie Warteschlange problemlos bilden.  Zuddem haben wir einen der Design Pattern \textbf{singleton} ausgewählt und umgesetzt. Durch diese zaubere Lösung könnten wir \textbf{SQL-Injektion} vermieden werden. Wir hatten außerdem die Möglichkeit durch die talentierte \textbf{Balzert}, unterschiedliche Analyse Muster und deren UML-repräsentation zu visualisieren. Da haben wir mögliche Problem ein der Muster sowie dessen Lösung untersuchit. Damit fehlerhafte Java-Formatierung nicht gepusht werden dürfen mussten wir einen Skript programmieren, den uns hilft effizient und problemlos Fehler zu schneller entdecken und sie abzulösen. Zur Visualierung der Auslastung unserer Anwendung haben wir ein Lasttest entwickelt, das die Anmeldung, Ausleihe und Abmeldung simuliert. Dadurch konnte das Verhalten der Webanwendung unter paralleler Belastung analysiert. 
